%% ieee-paper.tex
%% Template IEEE conference paper
%% Engine: pdflatex hoặc xelatex
%% Class: IEEEtran (download từ https://www.ieee.org/conferences/publishing/templates.html)
%%
%% Cách compile (tiếng Anh, pdflatex):
%%   pdflatex ieee-paper.tex
%%   biber ieee-paper
%%   pdflatex ieee-paper.tex
%%   pdflatex ieee-paper.tex

\documentclass[conference]{IEEEtran}

%% IEEEtran đã có sẵn nhiều cấu hình — không cần geometry

%% ================================================================
%% PACKAGES
%% ================================================================
\usepackage[T1]{fontenc}
\usepackage[utf8]{inputenc}

\usepackage{graphicx}
\usepackage{float}
\usepackage{booktabs}

%% Math
\usepackage{amsmath}
\usepackage{amssymb}

%% Cite với IEEEtran thường dùng biblatex ieee style
\usepackage[backend=biber, style=ieee, doi=true]{biblatex}
\addbibresource{sample.bib}

%% Hyperlinks (để ở cuối package list)
\usepackage{hyperref}
\hypersetup{
  colorlinks=false,    % IEEE thường không dùng colored links
  hidelinks=true
}

%% ================================================================
%% TITLE, AUTHORS, AFFILIATIONS
%% ================================================================
\title{Deep Learning for Vietnamese Text Classification:\\
       A Transformer-Based Approach}

\author{
  \IEEEauthorblockN{Van An Nguyen\IEEEauthorrefmark{1},
                    Thi Binh Tran\IEEEauthorrefmark{2}}
  \IEEEauthorblockA{\IEEEauthorrefmark{1}%
    School of Information and Communication Technology\\
    Hanoi University of Science and Technology\\
    Hanoi, Vietnam\\
    Email: an.nguyen@hust.edu.vn}
  \IEEEauthorblockA{\IEEEauthorrefmark{2}%
    Faculty of Computer Science\\
    Vietnam National University\\
    Hanoi, Vietnam\\
    Email: binh.tran@vnu.edu.vn}
}

%% ================================================================
%% BEGIN DOCUMENT
%% ================================================================
\begin{document}

\maketitle

%% ================================================================
%% ABSTRACT
%% ================================================================
\begin{abstract}
This paper presents a Transformer-based approach for Vietnamese text
classification. We fine-tune PhoBERT on a dataset of 50,000 Vietnamese
news articles spanning 10 categories. Our model achieves 92.7\% accuracy
on the test set, outperforming SVM (78.3\%) and LSTM (85.2\%) baselines
by a significant margin. Experiments demonstrate that pre-trained language
models substantially improve classification performance for low-resource
Vietnamese text.
\end{abstract}

%% ================================================================
%% KEYWORDS (IEEE format)
%% ================================================================
\begin{IEEEkeywords}
deep learning, text classification, Vietnamese NLP, Transformer,
BERT, fine-tuning
\end{IEEEkeywords}

%% ================================================================
%% MAIN BODY — IEEE uses 2-column layout
%% ================================================================

%% -------- I. INTRODUCTION --------
\section{Introduction}

Text classification is a fundamental task in natural language processing
(NLP) with applications ranging from spam detection to sentiment analysis
\cite{Goodfellow2016}. While significant progress has been achieved for
English, Vietnamese text classification remains challenging due to the
language's tonal nature and limited labeled resources \cite{Le2024bert}.

Recent advances in pre-trained language models, particularly BERT
\cite{Vaswani2017attention}, have shown remarkable performance across
various NLP tasks. However, direct application of English-trained models
to Vietnamese yields suboptimal results.

In this paper, we make the following contributions:
\begin{itemize}
  \item A fine-tuned PhoBERT model for Vietnamese text classification
  \item Comprehensive evaluation on VNews-50K, a new benchmark dataset
  \item Analysis of model behavior across text length and topic categories
\end{itemize}

%% -------- II. RELATED WORK --------
\section{Related Work}

\subsection{Text Classification Methods}

Traditional approaches to text classification rely on hand-crafted features
combined with classifiers such as Support Vector Machines (SVM) or Naive
Bayes \cite{Goodfellow2016}. While effective for well-defined feature spaces,
these methods require significant domain expertise.

Neural network-based approaches, including Convolutional Neural Networks
(CNN) and Long Short-Term Memory (LSTM) networks, learn representations
automatically from data \cite{LeCun2015deep}. However, they are trained
from scratch and lack transfer learning capabilities.

\subsection{Pre-trained Language Models}

The Transformer architecture \cite{Vaswani2017attention} revolutionized
NLP by enabling efficient attention over entire sequences. BERT and its
variants leverage large-scale pre-training on unlabeled text, enabling
effective transfer learning with minimal task-specific data.

For Vietnamese, PhoBERT \cite{Le2024bert} provides a domain-appropriate
pre-trained model trained on 20GB of Vietnamese text, demonstrating
strong performance on multiple downstream tasks.

%% -------- III. METHODOLOGY --------
\section{Methodology}

\subsection{Dataset}

We construct VNews-50K, a balanced dataset of 50,000 Vietnamese news
articles from three major outlets: VnExpress, Tuổi Trẻ, and Thanh Niên.
Articles are categorized into 10 topics with 5,000 samples each.
The dataset is split 70/15/15 for train/validation/test.

\subsection{Model Architecture}

Our approach fine-tunes PhoBERT-base \cite{Le2024bert} for sequence
classification. A linear classification head is added on top of the
\texttt{[CLS]} token representation:

\begin{equation}
  \label{eq:classifier}
  \hat{y} = \text{softmax}(\mathbf{W} \cdot h_{\text{CLS}} + \mathbf{b})
\end{equation}

where $h_{\text{CLS}} \in \mathbb{R}^{768}$ is the contextualized
representation of the \texttt{[CLS]} token, $\mathbf{W} \in \mathbb{R}^{10 \times 768}$,
and $\mathbf{b} \in \mathbb{R}^{10}$.

\subsection{Training Details}

The model is trained with the following hyperparameters:
\begin{itemize}
  \item Optimizer: AdamW \cite{Pytorch2024} with learning rate $2 \times 10^{-5}$
  \item Batch size: 32, max sequence length: 256 tokens
  \item Epochs: 10 with early stopping (patience = 3)
  \item Hardware: NVIDIA A100 40GB GPU
\end{itemize}

%% -------- IV. EXPERIMENTS --------
\section{Experiments and Results}

\subsection{Baselines}

We compare against four baselines:
(1) TF-IDF + SVM, (2) Word2Vec + CNN, (3) Word2Vec + BiLSTM,
and (4) mBERT (multilingual BERT).

\subsection{Main Results}

Table~\ref{tab:main_results} presents the classification results on
the test set. Our PhoBERT-based model achieves the best performance
across all metrics.

\begin{table}[htbp]
  \centering
  \caption{Performance Comparison on VNews-50K Test Set}
  \label{tab:main_results}
  \begin{tabular}{lcccc}
    \toprule
    Model      & Acc.   & Prec. & Rec.  & F1    \\
    \midrule
    TF-IDF+SVM & 78.3\% & 0.781 & 0.783 & 0.779 \\
    CNN        & 83.1\% & 0.829 & 0.831 & 0.828 \\
    BiLSTM     & 85.2\% & 0.851 & 0.852 & 0.849 \\
    mBERT      & 88.9\% & 0.887 & 0.889 & 0.886 \\
    \textbf{PhoBERT (Ours)} & \textbf{92.7\%} & \textbf{0.926} & \textbf{0.927} & \textbf{0.923} \\
    \bottomrule
  \end{tabular}
\end{table}

\subsection{Analysis}

\textbf{Effect of text length.} Our model maintains robust performance
for medium-length texts (100--300 tokens) but degrades for very short
texts ($<$50 tokens), achieving 85.4\% F1 in that range.

\textbf{Per-category performance.} The model performs best on
``Technology'' (F1 = 0.961) due to distinctive vocabulary, and
worst on ``Politics'' (F1 = 0.889) due to lexical overlap with
``Economics'' and ``Law''.

Fig.~\ref{fig:confusion} shows the confusion matrix, highlighting
the primary source of errors between related categories.

\begin{figure}[htbp]
  \centering
  %% Uncomment when figure file is available:
  %% \includegraphics[width=\columnwidth]{confusion_matrix.png}
  \fbox{\parbox{0.85\columnwidth}{\centering [Confusion Matrix Figure]\\
  \small Replace with actual \texttt{confusion\_matrix.png}}}
  \caption{Confusion matrix on VNews-50K test set (10 categories).
           Darker cells indicate higher confusion between categories.}
  \label{fig:confusion}
\end{figure}

%% -------- V. CONCLUSION --------
\section{Conclusion}

We presented a Transformer-based approach for Vietnamese text classification
that achieves state-of-the-art performance on VNews-50K. Fine-tuning
PhoBERT yields a 7.5-point improvement in accuracy over BiLSTM baseline.
Future work will explore data augmentation for short texts and
zero-shot generalization to new categories.

%% ================================================================
%% ACKNOWLEDGMENT (IEEE convention)
%% ================================================================
\section*{Acknowledgment}

The authors thank the reviewers for their insightful comments.
This work was supported by Vietnam National Foundation for Science
and Technology Development (NAFOSTED) under grant no.\ 102.01-2024.xx.

%% ================================================================
%% BIBLIOGRAPHY
%% ================================================================
\printbibliography

\end{document}
